% GENERATED FILE. Edit canonical lessons, then run npm run generate:books.
% canonical-source-hash: fnv1a64:13373288423967df
% canonical-lessons: SA-C521-vadakah, SA-C521-tantri, SA-C521-damaruh, SA-C521-murajah, SA-C521-nadah

\chapter{Drums and Strings}
\label{ch:sa-drums-and-strings}
\hlchaptermodality{521}

\begin{chapteropening}
\textbf{By the end of this chapter:} \emph{I can name instruments and say who plays them.}

Complete the last lesson of chapter 521: I can name instruments and say who plays them.
\end{chapteropening}

\section[\texorpdfstring{\sk{वादकः} (vādakaḥ)}{वादकः (vādakaḥ)}]{\sk{वादकः} (vādakaḥ) — a musician}

\label{lesson:SA-C521-vadakah}



\begin{quote}
Before the new one: say the Sanskrit for singing, then the Sanskrit for plays an instrument.
\end{quote}

\subsection*{You'll want to know: \sk{वादकः}}
\textbf{\sk{वादकः}} — \emph{vādakaḥ} — \textquotedblleft{}a musician\textquotedblright{}.

\begin{quote}
\sk{वादकः} \sk{मृदंगं} \sk{वादयति।}
\end{quote}

\begin{quote}
\emph{vādakaḥ mṛdaṅgaṃ vādayati.}
\end{quote}

\begin{quote}
The musician plays the mridanga.
\end{quote}

A masculine noun.

\begin{cousinweb}
\textbf{\sk{वादकः}} (\emph{vādakaḥ}) comes from the causative of the root √vad, \textquotedblleft{}to make sound\textquotedblright{}, with the agent ending \emph{-aka}; MW gives \emph{vādaka} as a musician (MW).
\end{cousinweb}

\begin{grammarlens}[title={What you've built}]
One more word for describing what you lived and felt.
\end{grammarlens}

\subsection*{Your turn}
\begin{itemize}
\raggedright
  \item \emph{Say it:} \emph{vādakaḥ}
  \item \emph{Say it:} \emph{vādakaḥ}, once more
  \item \emph{Say it:} \emph{vādakaḥ mṛdaṅgaṃ vādayati.} — The musician plays the mridanga.
  \item \emph{Recall:} say the Sanskrit for singing, then the Sanskrit for plays an instrument, then say \emph{vādakaḥ} again
\end{itemize}

\begin{tcolorbox}[breakable,colback=teal!4,colframe=teal!35!black,title={Before you move on}]
What does \sk{वादकः} mean? (\textquotedblleft{}A musician\textquotedblright{}.) Say it once more.
\end{tcolorbox}
\section[\texorpdfstring{\sk{तन्त्री} (tantrī)}{तन्त्री (tantrī)}]{\sk{तन्त्री} (tantrī) — a string of an instrument}

\label{lesson:SA-C521-tantri}



\begin{quote}
Before the new one: say the Sanskrit for plays an instrument, then the Sanskrit for a musician.
\end{quote}

\subsection*{You'll want to know: \sk{तन्त्री}}
\textbf{\sk{तन्त्री}} — \emph{tantrī} — \textquotedblleft{}a string of an instrument\textquotedblright{}.

\begin{quote}
\sk{वीणायाः} \sk{एका} \sk{तन्त्री} \sk{जीर्णा} \sk{अस्ति।}
\end{quote}

\begin{quote}
\emph{vīṇāyāḥ ekā tantrī jīrṇā asti.}
\end{quote}

\begin{quote}
One string of the veena is worn out.
\end{quote}

A feminine noun.

\begin{cousinweb}
\textbf{\sk{तन्त्री}} (\emph{tantrī}) comes from the root √tan, \textquotedblleft{}to stretch\textquotedblright{}, like \emph{tantu}, \textquotedblleft{}a thread\textquotedblright{}; MW gives \emph{tantrī} as the wire or string of a lute (MW).
\end{cousinweb}

\begin{grammarlens}[title={What you've built}]
One more word for describing what you lived and felt.
\end{grammarlens}

\subsection*{Your turn}
\begin{itemize}
\raggedright
  \item \emph{Say it:} \emph{tantrī}
  \item \emph{Say it:} \emph{tantrī}, once more
  \item \emph{Say it:} \emph{vīṇāyāḥ ekā tantrī jīrṇā asti.} — One string of the veena is worn out.
  \item \emph{Recall:} say the Sanskrit for plays an instrument, then the Sanskrit for a musician, then say \emph{tantrī} again
\end{itemize}

\begin{tcolorbox}[breakable,colback=teal!4,colframe=teal!35!black,title={Before you move on}]
What does \sk{तन्त्री} mean? (\textquotedblleft{}A string of an instrument\textquotedblright{}.) Say it once more.
\end{tcolorbox}
\section[\texorpdfstring{\sk{डमरुः} (ḍamaruḥ)}{डमरुः (ḍamaruḥ)}]{\sk{डमरुः} (ḍamaruḥ) — a hand drum}

\label{lesson:SA-C521-damaruh}



\begin{quote}
Before the new one: say the Sanskrit for a musician, then the Sanskrit for a string of an instrument.
\end{quote}

\subsection*{You'll want to know: \sk{डमरुः}}
\textbf{\sk{डमरुः}} — \emph{ḍamaruḥ} — \textquotedblleft{}a hand drum\textquotedblright{}.

\begin{quote}
\sk{बालः} \sk{डमरुं} \sk{वादयति।}
\end{quote}

\begin{quote}
\emph{bālaḥ ḍamaruṃ vādayati.}
\end{quote}

\begin{quote}
The child plays a hand drum.
\end{quote}

A masculine noun.

\begin{cousinweb}
\textbf{\sk{डमरुः}} (\emph{ḍamaruḥ}) is the word MW records for a sacred drum shaped like an hourglass, the drum of Śiva; MW gives no derivation, so its origin is uncertain (MW).
\end{cousinweb}

\begin{grammarlens}[title={What you've built}]
One more word for describing what you lived and felt.
\end{grammarlens}

\subsection*{Your turn}
\begin{itemize}
\raggedright
  \item \emph{Say it:} \emph{ḍamaruḥ}
  \item \emph{Say it:} \emph{ḍamaruḥ}, once more
  \item \emph{Say it:} \emph{bālaḥ ḍamaruṃ vādayati.} — The child plays a hand drum.
  \item \emph{Recall:} say the Sanskrit for a musician, then the Sanskrit for a string of an instrument, then say \emph{ḍamaruḥ} again
\end{itemize}

\begin{tcolorbox}[breakable,colback=teal!4,colframe=teal!35!black,title={Before you move on}]
What does \sk{डमरुः} mean? (\textquotedblleft{}A hand drum\textquotedblright{}.) Say it once more.
\end{tcolorbox}
\section[\texorpdfstring{\sk{मुरजः} (murajaḥ)}{मुरजः (murajaḥ)}]{\sk{मुरजः} (murajaḥ) — a tambourine}

\label{lesson:SA-C521-murajah}



\begin{quote}
Before the new one: say the Sanskrit for a string of an instrument, then the Sanskrit for a hand drum.
\end{quote}

\subsection*{You'll want to know: \sk{मुरजः}}
\textbf{\sk{मुरजः}} — \emph{murajaḥ} — \textquotedblleft{}a tambourine\textquotedblright{}.

\begin{quote}
\sk{उत्सवे} \sk{जनाः} \sk{मुरजं} \sk{वादयन्ति।}
\end{quote}

\begin{quote}
\emph{utsave janāḥ murajaṃ vādayanti.}
\end{quote}

\begin{quote}
At the festival people play the tambourine.
\end{quote}

A masculine noun.

\begin{cousinweb}
\textbf{\sk{मुरजः}} (\emph{murajaḥ}) is the word MW records for a kind of drum, a tambourine; MW suggests \emph{mura} and \emph{ja}, \textquotedblleft{}born\textquotedblright{}, with a query, so its origin is uncertain (MW).
\end{cousinweb}

\begin{grammarlens}[title={What you've built}]
One more word for describing what you lived and felt.
\end{grammarlens}

\subsection*{Your turn}
\begin{itemize}
\raggedright
  \item \emph{Say it:} \emph{murajaḥ}
  \item \emph{Say it:} \emph{murajaḥ}, once more
  \item \emph{Say it:} \emph{utsave janāḥ murajaṃ vādayanti.} — At the festival people play the tambourine.
  \item \emph{Recall:} say the Sanskrit for a string of an instrument, then the Sanskrit for a hand drum, then say \emph{murajaḥ} again
\end{itemize}

\begin{tcolorbox}[breakable,colback=teal!4,colframe=teal!35!black,title={Before you move on}]
What does \sk{मुरजः} mean? (\textquotedblleft{}A tambourine\textquotedblright{}.) Say it once more.
\end{tcolorbox}
\section[\texorpdfstring{\sk{नादः} (nādaḥ)}{नादः (nādaḥ)}]{\sk{नादः} (nādaḥ) — a resonant sound}

\label{lesson:SA-C521-nadah}



\begin{quote}
Before the new one: say the Sanskrit for a hand drum, then the Sanskrit for a tambourine.
\end{quote}

\subsection*{You'll want to know: \sk{नादः}}
\textbf{\sk{नादः}} — \emph{nādaḥ} — \textquotedblleft{}a resonant sound\textquotedblright{}.

\begin{quote}
\sk{मृदंगस्य} \sk{नादः} \sk{दूरं} \sk{गच्छति।}
\end{quote}

\begin{quote}
\emph{mṛdaṅgasya nādaḥ dūraṃ gacchati.}
\end{quote}

\begin{quote}
The sound of the drum carries far.
\end{quote}

A masculine noun.

\begin{cousinweb}
\textbf{\sk{नादः}} (\emph{nādaḥ}) comes from the root √nad, \textquotedblleft{}to sound, to roar\textquotedblright{}; MW gives \emph{nāda} as a loud sound, a roar, and any sound or tone (MW).
\end{cousinweb}

\begin{grammarlens}[title={What you've built}]
One more word for describing what you lived and felt.
\end{grammarlens}

\subsection*{Your turn}
\begin{itemize}
\raggedright
  \item \emph{Say it:} \emph{nādaḥ}
  \item \emph{Say it:} \emph{nādaḥ}, once more
  \item \emph{Say it:} \emph{mṛdaṅgasya nādaḥ dūraṃ gacchati.} — The sound of the drum carries far.
  \item \emph{Recall:} say the Sanskrit for a hand drum, then the Sanskrit for a tambourine, then say \emph{nādaḥ} again
\end{itemize}

\begin{tcolorbox}[breakable,colback=teal!4,colframe=teal!35!black,title={Before you move on}]
What does \sk{नादः} mean? (\textquotedblleft{}A resonant sound\textquotedblright{}.) Say it once more.
\end{tcolorbox}
